\section{Proofs for parametric support certificates}\label{sec:parametric-proofs}

\subsection{Proof of Theorem~\ref{thm:cover-con}}\label{sec:parametric-proofs-cover}

\begin{proof}
Fix $i$ and $p\ne q$ in $\mathcal D$, and put $\gamma(t)=q+t(p-q)$ for
$t\in[0,1]$; convexity keeps $\gamma(t)$ in $\mathcal D$. The sets
$S_s=\{t:\gamma(t)\in C_{i,s}\}$ are finite unions of closed intervals and
cover $[0,1]$. Let $\mathcal J$ be the finite collection of the intervals
of positive length among their components. The union of $\mathcal J$ is
closed, and its complement in $[0,1]$ is relatively open and contained in a
finite set, hence empty. Let $0=t_0<t_1<\cdots<t_q=1$ be the endpoints of
the intervals in $\mathcal J$. The midpoint of $(t_{r-1},t_r)$ lies in some
$[\alpha,\beta]\in\mathcal J$ with $[\alpha,\beta]\subseteq S_s$. Since
$\alpha$ and $\beta$ are among the $t$'s and none lies strictly between
$t_{r-1}$ and $t_r$, we get $[t_{r-1},t_r]\subseteq S_s$. On this interval
$F_i(\gamma(t))=h_{i,s}(\gamma(t))$, a continuously differentiable function
of $t$ whose derivative
$\sum_j\partial_jh_{i,s}(\gamma(t))(p_j-q_j)$ has absolute value at most
$\sum_jM_{ij}|p_j-q_j|$. The mean value theorem gives
\[
 |F_i(\gamma(t_r))-F_i(\gamma(t_{r-1}))|
 \le(t_r-t_{r-1})\sum_jM_{ij}|p_j-q_j|.
\]
Summing over $r$ proves $|F_i(p)-F_i(q)|\le\sum_jM_{ij}|p_j-q_j|$. For
$q\le p$, the two-sided bound gives $F(p)-F(q)\le M(p-q)$.
\end{proof}

\subsection{Proof of Proposition~\ref{prop:coverage-con}}\label{sec:parametric-proofs-coverage}

\begin{proof}
Each program is feasible, since any $\theta\in\mathcal D$ with a
sufficiently negative $t$ is feasible, and its optimal value is attained,
since it is the maximum over the compact set $\mathcal D$ of the continuous
function $\min\{1,H^r_{j_r}\theta-g^r_{j_r}\}$. If some
$\theta\in\mathcal D$ lies in no $C^r$, choose for each $r$ an inequality
of $C^r$ that $\theta$ violates strictly. The minimum of these violations
and $1$ is positive and gives $\tau_J>0$ for the selected tuple.
Conversely, a feasible point with $t>0$ violates strictly one inequality of
every $C^r$ and so lies outside the union.
\end{proof}
